\documentclass[11pt,a4paper]{article}
\usepackage[margin=.75in]{geometry}
\usepackage{xcolor}
\usepackage[T1]{fontenc}
\usepackage{lmodern,microtype,amsmath,amssymb,amsthm,graphicx,booktabs,enumitem,placeins}
\usepackage[colorlinks=true,linkcolor=blue!45!black,urlcolor=blue!45!black,citecolor=blue!45!black]{hyperref}
\pagestyle{empty}
\setlength{\parindent}{0pt}\setlength{\parskip}{3pt}
\setlength{\textfloatsep}{12pt}\setlength{\intextsep}{12pt}
\makeatletter
\setlength{\@fptop}{0pt}\setlength{\@fpsep}{12pt}\setlength{\@fpbot}{0pt plus 1fil}
\makeatother
\setcounter{tocdepth}{2}
\newtheorem{theorem}{Theorem}\newtheorem{lemma}{Lemma}
\newcommand{\E}{\mathbb E}\newcommand{\Prb}{\mathbb P}
\newcommand{\fig}[3]{\begin{figure}[!htbp]\centering\includegraphics[width=#2\linewidth]{../figures/#1.pdf}\caption{#3}\label{fig:#1}\end{figure}}
\hypersetup{pdftitle={COBRA: Confidence-Optimized Borda Ranking Algorithm},pdfauthor={Pranav Bhatnagar and Praveen Kumar S}}
\begin{document}
\hypersetup{pageanchor=false}
\begin{titlepage}\centering\vspace*{1.0cm}
{\Large IIT Madras\\DA5453: Learning from Human Preferences\par}
\vspace{1.2cm}{\Huge\bfseries COBRA\par}\vspace{.4cm}
{\LARGE Confidence-Optimized Borda\\Ranking Algorithm\par}\vspace{.7cm}
{\Large An investigation of adaptive comparisons\\and confidence-based information reuse\par}
\vspace{1.2cm}
\begin{tabular}{ll}\toprule Team member&Roll number\\\midrule
Pranav Bhatnagar&ED23B037\\Praveen Kumar S&ED23B038\\\bottomrule\end{tabular}
\vfill\begin{minipage}{.88\linewidth}
\textbf{Research question.} Can adaptive pair selection and statistically valid reuse of duel outcomes identify the Borda winner with fewer total comparisons than uniform-opponent Borda learning?

\textbf{Finding.} The answer depends on the instance. Reuse can offset exploration costs on small problems; conservative pairwise intervals and extra adaptive exploration can make the method more expensive as the number of alternatives grows. Our contribution is a controlled investigation with a valid certificate and a safeguard guarantee, rather than a claim of universal improvement.
\end{minipage}\vspace*{1cm}\end{titlepage}
\clearpage\hypersetup{pageanchor=true}\pagenumbering{roman}
{\small\tableofcontents}\clearpage\pagenumbering{arabic}

\section{Introduction and research question}
Pairwise feedback is natural when people can compare two alternatives more reliably than assign absolute scores. Examples include judging two performances, choosing a preferred search result, and comparing candidates on a ballot. A learner must decide which comparison to request next: feedback can be costly even when computation is cheap. Our starting paper, RUCB~\cite{zoghi}, learns under a strict Condorcet-winner assumption. Such an alternative beats every other alternative with probability greater than one half. Cycles can invalidate that assumption, and even when a Condorcet winner exists, another aggregation rule may prefer a different alternative.

This issue matters in real preference aggregation. ERS ballots come from elections held by non-profit organizations, trade unions, and professional bodies~\cite{ers}; our full 26-candidate profile contains 13 strict three-cycles and no strict Condorcet winner. Skate rankings aggregate judges' assessments of real competitions~\cite{skate}; our full profile has local cycles, although it does have a Condorcet winner. Thus disagreement need not produce a single alternative that beats everyone, and cycles alone do not prove that none exists. A Condorcet-only learner cannot supply its intended winner when that winner is absent.

We choose Borda as the objective from the first duel: it averages winning probabilities against all other alternatives. A maximizer always exists, although ties are possible. This is a change of aggregation criterion, not a claim that Borda is always preferable to Condorcet. It avoids a diagnostic stage that can be expensive near pairwise ties. COBRA then asks whether shared evidence and adaptive comparisons improve on straightforward Borda learning.

The central question is whether \emph{adaptive pair selection and shared confidence information reduce the total number of comparisons needed for Borda identification}. A simple baseline draws a uniform opponent for each designated arm; its binary outcome is already an unbiased Borda sample. COBRA retains each duel in a reciprocal pairwise matrix, constructs valid score intervals, and devotes one quarter of exploration blocks to uncertain comparisons involving a current leader or rival. Three quarters retain the baseline sampling mechanism to guarantee progress. This is a deliberately modest heuristic with a conservative proof.

Our main finding is that adaptive confidence reuse offers modest improvements on selected small instances, but its overhead can outweigh its benefits as $K$ grows. Across @@RUNS@@ primary and @@HRUNS@@ horizon-aware simulations, COBRA saves 6.8\% on the three-arm cycle (95\% paired bootstrap interval 2.5--10.8\%), alongside \emph{31.6\% and 36.7\% additional cost} on the planted $K=5$ and $K=10$ instances. The planted three-arm saving is @@S3@@\%, with an interval including zero. Both methods make correct empirical recommendations in 30/30 runs on each real profile, but certify in 0/30: these are not statistically certified identifications. The added weight-sensitivity experiment tests whether the heuristic's design choice changes this picture.

\section{Problem formulation and motivating examples}
There are $K\ge2$ alternatives, $\mathcal A=\{1,\dots,K\}$. A duel between distinct alternatives $i,j$ returns $Y=1$ if $i$ wins, with $Y\sim\operatorname{Bernoulli}(p_{ij})$. We assume stationary probabilities, $p_{ij}+p_{ji}=1$, $p_{ii}=1/2$, and independent draws conditional on the selected pair. Pair selection may depend on all past feedback. Neither transitivity nor a parametric score model is assumed. The policy receives outcomes and its own counts; the true matrix is available only to the simulator and evaluator.

Define
\begin{equation}\label{eq:borda}
b_i=\frac1{K-1}\sum_{j\ne i}p_{ij},\qquad
i^\star\in\arg\max_i b_i,\qquad\Delta_i=b_{i^\star}-b_i.
\end{equation}
For the termination theorem assume a unique maximizer and write $\Delta_B=\min_{i\ne i^\star}\Delta_i>0$. The implementation is not given these scores or gaps. Exact ties remain valid inputs, but a strict-winner termination guarantee is inapplicable. At a budget cap we return an explicitly uncertified empirical recommendation if no certificate has appeared.

\textbf{Different winners.} Consider
\[
P=\begin{pmatrix}.5&.51&.51\\.49&.5&.90\\.49&.10&.5\end{pmatrix},\qquad
b=(.51,.695,.295).
\]
Alternative $A$ is the strict Condorcet winner, but $B$ is the unique Borda winner. COBRA must choose $B$. Calling that outcome a failure to preserve Condorcet priority would evaluate the wrong objective. Borda rewards average performance, and this choice of aggregation rule is substantive rather than universally correct.

\textbf{A unique winner inside a cycle.} With $p_{12}=.866\overline6$, $p_{23}=.8$, and $p_{31}=.733\overline3$, the preferences cycle as $1\succ2\succ3\succ1$. The Borda scores are $(.566\overline6,.466\overline6,.466\overline6)$, with a unique winner and gap $.1$. A symmetric cycle with all three directed probabilities $.8$ instead has three equal Borda scores. Existence of cycles and uniqueness of the Borda winner are separate issues.

\textbf{Two costs.} Identification cost is the number $\tau$ of duels needed for a statistical certificate. Strong Borda regret measures the quality of \emph{both} alternatives requested:
\begin{equation}\label{eq:regret}
r_t=b_{i^\star}-\frac{b_{I_t}+b_{J_t}}2,\qquad R_T=\sum_{t=1}^T r_t.
\end{equation}
It lies in $[0,1]$ per round. After certification the policy uses self-comparisons $(\hat i,\hat i)$, which incur zero regret if $\hat i$ is correct; these are exploitation actions, not informative Bernoulli samples. Fewer exploration duels need not imply less regret when different pairs have different losses. All reported regret uses~\eqref{eq:regret}, including on matrices that possess a Condorcet winner.

\section{Related work and scope of the contribution}
RUCB~\cite{zoghi} motivates this project but targets Condorcet; its regret theorem does not transfer. Jamieson et al.~\cite{jamieson} study Borda identification, the uniform-opponent reduction, and Successive Elimination with Comparison Sparsity (SECS). SECS compares active alternatives with a random reference and adds a sparsity-based elimination test. Their lower bound, for probabilities in $[3/8,5/8]$ and $K\ge4$, has inverse-square Borda-gap dependence. COBRA neither implements SECS nor inherits its structural assumptions or theorem.

Heckel et al.~\cite{heckel} study active Borda-type ranking without structural assumptions. Shah and Wainwright~\cite{shah} establish the strength of simple Borda counting under random comparisons, motivating a serious uniform baseline. Saha et al.~\cite{saha} analyze adversarial Borda regret, including a fixed-gap setting. Their results are not imported into our stochastic finite-horizon proof.

Confidence intervals, Hoeffding bounds, intersection, and certificates are standard. We investigate their combination with a leader/rival heuristic and scheduled uniform exploration. Adaptive Borda learning and informative-opponent selection are established ideas; no publication-level novelty is claimed. Our contribution is a matched empirical investigation and a self-contained safeguard proof.

\section{Method: COBRA}
\subsection{Shared evidence and two valid estimators}
Store a win matrix $W$, where $w_{ij}$ counts outcomes in which $i$ beat $j$, and $N_{ij}=w_{ij}+w_{ji}$. Each physical duel contributes one win and one increment to the count of its unordered pair. Mirrored count entries represent the same observation; they are not two independent samples. For $n=N_{ij}>0$, set $\hat p_{ij}=w_{ij}/n$. Write $M=K(K-1)/2$ and split the error budget as $\delta_D=\delta_P=\delta/2$.

For $q$ objects define
\begin{equation}\label{eq:radius}
\rho_q(n;\eta)=\sqrt{\frac{\log(4qn^2/\eta)}{2n}}\quad(n\ge1).
\end{equation}
Each pair interval is $[L_{ij},U_{ij}]=[\hat p_{ij}-\rho_M,\hat p_{ij}+\rho_M]\cap[0,1]$. Unseen pairs have $[0,1]$; diagonal entries equal $1/2$. Reverse intervals satisfy $L_{ji}=1-U_{ij}$. The row intervals are
\begin{equation}\label{eq:rows}
\ell_i^P=\frac{\sum_{j\ne i}L_{ij}}{K-1},\qquad
u_i^P=\frac{\sum_{j\ne i}U_{ij}}{K-1}.
\end{equation}
Adaptive allocation does not invalidate these intervals: they concern the probability of a particular pair, not a mixture of nonuniform opponents.

In a \emph{safe} block, designate the least-sampled arm according to direct-sample counts $c_i$, breaking ties by index, and draw each opponent independently and uniformly from its other $K-1$ arms. Only the designated arm receives a direct Borda sample. Its sum $s_i$ and count $c_i$ yield direct intervals $[\ell_i^D,u_i^D]$ using~\eqref{eq:radius} with $q=K$, error budget $\delta_D$, and empirical mean $s_i/c_i$. Before its first sample this interval is $[0,1]$. Adaptive-block outcomes do not enter $s_i,c_i$.

Combine the evidence by intersection:
\begin{equation}\label{eq:intersection}
\ell_i=\max(\ell_i^D,\ell_i^P),\qquad u_i=\min(u_i^D,u_i^P).
\end{equation}
These estimators can be correlated; simultaneous coverage, rather than independence between estimators, justifies intersection. Naively averaging all adaptively chosen wins would target the chosen-opponent distribution instead of~\eqref{eq:borda}. COBRA never makes that substitution.

\subsection{Adaptive selection, safeguard, and certificate}
Use blocks of $B=128$ duels, freezing the choice rule within each block. Three safe blocks are followed by one adaptive block. In an adaptive block, compute row means from $\hat P$, substituting $1/2$ for unseen probabilities. Let $a$ maximize this empirical score. Let $b\ne a$ maximize the pairwise row upper bound $u_b^P$. The pair $a,b$ is not necessarily the next duel: the difference satisfies
\begin{equation}\label{eq:diff}
(K-1)(b_a-b_b)=2p_{ab}-1+\sum_{j\notin\{a,b\}}(p_{aj}-p_{bj}).
\end{equation}
Equation~\eqref{eq:diff} explains the weights: $p_{ab}$ has coefficient 2 because the duel enters the two scores with opposite signs, whereas each other incident probability has coefficient $+1$ or $-1$. Assign the corresponding absolute weights 2 and 1 (zero elsewhere), and maximize weight times width $U_{ij}-L_{ij}$, with uniform random tie-breaking. The common factor $1/(K-1)$ cancels. These coefficients motivate the heuristic; they do not prove its allocation optimal. It ignores the samples needed to shrink a width and the possibility that the leader changes. We test alternative leader--rival weights 1 and 4 while retaining all other incident weights at 1.

Stop at the first inspected block boundary with
\begin{equation}\label{eq:certificate}
\ell_i>\max_{j\ne i}u_j.
\end{equation}
Return the unique certified arm. At the cap return the largest empirical score, explicitly marked uncertified. No arm is permanently eliminated, so all opponents remain available for safe sampling. A nearly tied pair is not tested merely to decide whether it crosses one half.

\begin{center}\fbox{\begin{minipage}{.94\linewidth}
\textbf{Algorithm 1: Confidence-Optimized Borda Ranking Algorithm}\par
\textbf{Input:} $K$, error level $\delta$, block size $B$, duel cap $H$; no matrix or gap.\par
Initialize $W,N,s,c=0$ and block index $q=0$.\par
While fewer than $H$ duels have been used:
\begin{enumerate}[nosep,leftmargin=1.5em]
\item If $q\bmod4\ne3$, choose $i=\arg\min c_i$ and run a safe block with fresh uniform opponents. Otherwise choose $(i,j)$ by~\eqref{eq:diff}'s weighted-width rule and run a repeated-pair block.
\item Update the reciprocal pair matrix once per duel. Update $s_i,c_i$ only for designated safe observations.
\item Compute~\eqref{eq:radius}--\eqref{eq:intersection}. If~\eqref{eq:certificate} holds, return the certified arm. Otherwise increment $q$.
\end{enumerate}
At exhaustion return an empirical recommendation with its uncertified status. For regret evaluation, self-duel a certified arm for all remaining rounds.
\end{minipage}}\end{center}
\fig{schematic}{.95}{The two sampling branches feed the same pair matrix. Only the safe branch feeds the direct Borda estimator. Both estimators support a common stopping certificate. The evaluator alone knows the ground-truth matrix and Borda scores.}

\textbf{Worked updates.} Take the three-arm cycle from Section 2 and use $B=1$ only to make the bookkeeping visible. Initially all intervals are $[0,1]$. In iteration 1 designate arm 1; suppose the random opponent is 2 and arm 1 wins. Then $w_{12}=1,w_{21}=0,N_{12}=N_{21}=1$, and $(s_1,c_1)=(1,1)$. In iteration 2, arm 2 is the least-sampled arm; suppose its opponent is 3 and it loses. Then $w_{32}=1,N_{23}=N_{32}=1$, and $(s_2,c_2)=(0,1)$. No direct sample is added for arm 3. At this confidence level both observed-pair intervals are still clipped to $[0,1]$, so neither iteration certifies a winner despite extreme empirical means. Iteration 3 designates arm 3. Iteration 4 applies the adaptive rule, which gives the leader--rival edge twice the weight of their other incident edges. Production experiments use $B=128$ and integer aggregate counts; the bookkeeping and certificate are identical.

\section{Theoretical analysis}
\subsection{A. Uniform-opponent concentration}
\begin{theorem}[Fixed-count baseline]
Fix $\delta\in(0,1)$ and a unique Borda winner. If every arm receives
\[
n_B=\left\lceil\frac8{\Delta_B^2}\log\frac{2K}{\delta}\right\rceil
\]
designated uniform-opponent samples, the empirical largest score identifies the winner with probability at least $1-\delta$, using $Kn_B$ duels.
\end{theorem}
\begin{proof}
For a designated arm $i$, conditional on the history its next fresh opponent is uniform and independent, giving $\Prb(Y=1)=\sum_{j\ne i}p_{ij}/(K-1)=b_i$. Independent fresh draws give an i.i.d. Bernoulli sequence for that arm. Hoeffding gives $\Prb(|\hat b_i-b_i|>\epsilon)\le2e^{-2n\epsilon^2}$. With $\epsilon=\Delta_B/4$ and the displayed $n_B$, a union bound across $K$ arms gives simultaneous accuracy with failure probability at most $\delta$. On that event $\hat b_{i^\star}-\hat b_j\ge\Delta_j-\Delta_B/2\ge\Delta_B/2>0$ for every competitor.
\end{proof}
This is an analysis budget depending on an unknown gap. The implementation does not choose its sample count from $n_B$. Moreover, this fixed-count statement alone does not validate repeatedly inspecting fixed-$n$ intervals. For the actual sequential stopping rule we use the next lemma.

\subsection{B. Simultaneous confidence and adaptive reuse}
\begin{lemma}[Count-uniform coverage]\label{lem:coverage}
For $q$ Bernoulli streams with fixed respective means, intervals using~\eqref{eq:radius} simultaneously cover all means at every positive local sample count, with failure probability at most $\eta\pi^2/12<\eta$. Streams need not be mutually independent. An adaptive rule may decide which stream to reveal next, provided each revealed draw has its stream's stationary independent-draw law.
\end{lemma}
\begin{proof}
For one stream and a fixed count $n$, Hoeffding bounds its two-sided failure probability by
\[
2e^{-2n\rho_q(n;\eta)^2}=\frac{\eta}{2qn^2}.
\]
Summing over $q$ streams and all $n\ge1$ gives $\eta\sum_{n\ge1}1/(2n^2)=\eta\pi^2/12$. Equivalently, imagine an infinite sequence of outcomes for each unordered pair generated in advance. An adaptive choice reveals the next unused outcome from one sequence. The union event already covers every possible count, including the random counts at stopping. No conditioning on a chosen random count is substituted into a fixed-count probability bound. Clipping to $[0,1]$ cannot remove the true mean.
\end{proof}
Apply the lemma to $M$ unordered-pair streams with $\eta=\delta_P$ and to $K$ direct Borda streams with $\eta=\delta_D$. Safe designated outcomes have the required law even when the designated arm depends on past observations. The two events can be dependent, but their failure probabilities add to at most $\delta\pi^2/12<\delta$. Reciprocity supplies the reverse bounds without allocating another independent confidence budget. Taking averages proves $b_i\in[\ell_i^P,u_i^P]$ for every arm, and intersection with the direct intervals preserves coverage. The directional ablation instead treats $K(K-1)$ ordered streams separately and pays that larger union-bound factor.

\begin{theorem}[Valid certificate]
With probability at least $1-\delta$, every certificate emitted by COBRA, at any inspected time, names the unique true Borda winner.
\end{theorem}
\begin{proof}
On simultaneous coverage,~\eqref{eq:certificate} implies $b_i\ge\ell_i>u_j\ge b_j$ for each $j\ne i$. Thus a wrong or nonunique winner cannot pass the certificate on that event. This conclusion concerns correctness, not guaranteed termination at a finite cap.
\end{proof}

\subsection{C. Safeguarded sample complexity and regret}
Let $\bar n$ be any integer such that $\rho_K(n;\delta_D)\le\Delta_B/8$ for every $n\ge\bar n$. It is finite and is used only in analysis. With equal complete blocks, after $S$ safe blocks the designated counts differ by at most $B$: the next safe block always goes to a current minimum. Therefore $S B\ge K(\bar n+B)$ is sufficient for every $c_i\ge\bar n$. If at most $J$ otherwise-full blocks are shortened to hit logging checkpoints, add $J$ safe blocks conservatively. Define
\begin{equation}\label{eq:N0}
S_0=\left\lceil\frac{K(\bar n+B)}B\right\rceil+J,\qquad
N_0=4B\left\lceil\frac{S_0}{3}\right\rceil.
\end{equation}
An untruncated run has $J=0$; the experiments force at most three interior checkpoints, so $J\le3$. Every four scheduled blocks contain three safe blocks; an incomplete final block at the cap cannot increase this bound on work before certification in an uncapped run.

\textbf{Implementation audit.} After a shortened block the runner increments its index once, continuing the $q\bmod4$ schedule without restarting it. At most $J$ safe blocks each lose fewer than $B$ samples. After $S$ safe blocks, the total is at least $(S-J)B$. Least-count designation keeps the largest and smallest counts within $B$; the minimum is therefore at least $(S-J)B/K-B$, proving the choice of $S_0$. Tests exercise shortened blocks, the $N_0$ count bound over 108 parameter combinations, the confidence split, intersection, and strict certificate. The terminal cap is separate from logging cuts; a run that ends there without a certificate remains capped.

\begin{theorem}[Conservative safeguard guarantee]\label{thm:sample}
Under the model assumptions and $\Delta_B>0$, on the simultaneous confidence event COBRA certifies the true winner by $N_0$ duels. Thus $\Prb(\tau>N_0)\le\delta$, where an uncapped process is understood. The adaptive choice rule can be replaced arbitrarily without invalidating this guarantee, provided the three-safe/one-adaptive schedule and certificate are retained.
\end{theorem}
\begin{proof}
At $c_i\ge\bar n$, both the empirical error and direct interval radius are at most $\Delta_B/8$ on the good event. Hence $\ell_{i^\star}^D\ge b_{i^\star}-\Delta_B/4$ and $u_j^D\le b_j+\Delta_B/4$. Their separation is at least $\Delta_j-\Delta_B/2\ge\Delta_B/2>0$. Intersection only raises lower bounds and lowers upper bounds. Once the safe counts reach the stated threshold, the common certificate must therefore hold. Balancing safe counts and the block schedule gives~\eqref{eq:N0}. Adaptive observations can help earlier but are not needed for this worst-case argument.
\end{proof}
For an explicit order, put $x=\Delta_B^{-2}\ge1$, $L=\log(256Kx/\delta_D)$, and $\bar n=\lceil512xL\rceil$. For $n=\bar n\le513xL$, direct substitution yields
\[
\log(4Kn^2/\delta_D)\le2L+2\log L+8.33\le4L,
\]
since $L\ge\log512$. Consequently $\rho_K(n;\delta_D)^2\le1/(256x)<\Delta_B^2/64$. The radius decreases for $n\ge1$ here: differentiating its squared expression gives a derivative with sign $2-\log(4Kn^2/\delta_D)<0$. This proves the required condition for all larger $n$ and gives
\begin{equation}\label{eq:sampleorder}
N_0=O\!\left(\frac K{\Delta_B^2}\log\frac K{\delta\Delta_B}+KB+JB\right).
\end{equation}
This conservative constant is not a prediction of observed stopping time. Nor does the proof establish a better dependence for the adaptive heuristic. It covers the full implemented policy through its safeguard.

\begin{theorem}[Finite-horizon expected Borda regret]
Run COBRA at confidence $\delta$ until certification or horizon $T$, then self-duel any certified winner. Its regret satisfies
\[
\E R_T\le\min(T,N_0)+T\delta.
\]
With a supplied finite horizon $T\ge2$ and $\delta=T^{-2}$, fixed $B,J$ give
\[
\E R_T=O\!\left(\min\!\left\{T,\frac K{\Delta_B^2}\log\frac{KT}{\Delta_B}+KB+JB\right\}\right)+T^{-1}.
\]
\end{theorem}
\begin{proof}
On the good event, at most $\min(T,N_0)$ exploration rounds each cost at most one, followed by correct zero-regret exploitation. On its complement the total cost is at most $T$. Taking expectations proves the first statement; substituting $T^{-2}$ into~\eqref{eq:sampleorder} proves the second.
\end{proof}
For fixed positive gap and fixed $K$, this is logarithmic in $T$; its displayed dependence includes $K$, gap, and logarithmic factors. Writing simply $O(K\log T)$ would hide the crucial $\Delta_B^{-2}$ dependence and the $\log K$ term. A fixed-$\delta=.05$ identification run gives time-uniform correctness but the regret argument still has a $.05T$ failure term. An anytime confidence interval is not automatically an anytime logarithmic expected-regret theorem. The separate $T^{-2}$ simulations use a different error level for each supplied horizon.

\subsection{D. Computation, limitations of the bound}
Direct score maintenance needs $O(K)$ learner statistics and $O(1)$ work per observation, with $O(K)$ work for a full certificate check. COBRA stores $O(K^2)$ pair statistics, makes $O(1)$ pair updates per duel, and recomputes pair intervals, row bounds, and adaptive scores in $O(K^2)$ work per inspected block. Its exploration cost is therefore $O(\tau K^2/B+\tau)$ in the straightforward implementation. A safe simulation block aggregates opponents into at most $K-1$ integer counts; this reduces simulation overhead but not the charged number of duels. All timings include policy computation and simulation, not hypothetical human response latency.

The benchmark also retains the full comparison-count matrix for uniform baselines to measure allocation. This instrumentation uses $O(K^2)$ memory in our runner, although the baseline's decision rule itself requires only $O(K)$ memory. Our proof supplies a gap-dependent upper bound, not an information-theoretically optimal constant, an expected stopping-time bound on the bad event, or a claim of $\sqrt K$ scaling. The lack of adaptive-specific improvement is deliberate: formal correctness is retained even if the heuristic allocates poorly.

\section{Experimental methodology}
\subsection{Algorithms and fairness controls}
\textbf{Uniform Borda (fixed budget)} cycles through least-sampled designated arms, draws uniform opponents, and returns an uncertified recommendation at the cap. \textbf{Uniform + stop} uses the same rule but checks direct anytime intervals after each block, spending its whole error budget on direct scores. Comparing these isolates stopping rather than informative pair allocation. \textbf{COBRA} uses the full method above. \textbf{Random adaptive block} replaces only the one adaptive block by a uniformly random unordered pair, keeping the safeguard and confidence intersection unchanged. This is the main allocation control.

\textbf{Direct-only certificate} retains COBRA's shadow pair matrix solely for the identical adaptive selection rule, but ignores its row intervals for stopping; it uses the same $\delta/2$ direct budget as COBRA. Up to the earlier stopping time, these two variants have identical allocation and feedback under the same seed. Their difference isolates access to the pairwise certificate. \textbf{Forward evidence only} also keeps this same shadow allocation, but discards mirrored outcomes when constructing its certificate's ordered-pair intervals. This tests the value of reciprocal evidence, while paying the appropriate larger union-bound correction; it is not an independence claim. \textbf{No safeguard} makes every block adaptive and has a valid certificate but no proved sample bound. It is a secondary candidate, not the chosen algorithm. The safeguarded method is primary because its guarantee does not depend on heuristic exploration behavior; it was not selected after searching the performance table.

All modes receive the same matrix, cap, and error level, and use the same Borda objective. Each replication has separate selection and feedback random generators, initialized by seed $2026000+s$ and seed plus $10^6$, respectively. Identical-allocation ablations share feedback exactly until one stops. Other algorithms share starting seeds but need not consume the same pair streams, so we do not claim exact outcome coupling under different allocations. Batch size is 128 for every mode. Blocks are shortened at prescribed horizons, without observing outcomes to choose their lengths. Opponent multinomials followed by binomial pair counts simulate exactly the distribution of independent uniform-opponent draws under a block-frozen designation. These are integer comparisons, not fractional pseudo-observations.

\subsection{Experiments A--D: synthetic preference matrices}
\textbf{A: number of alternatives.} Use $K\in\{3,5,10,20,50\}$ and set $p_{0j}=.5+a$, $p_{j0}=.5-a$, all other distinct probabilities to $.5$, with $a=g(K-1)/K$ and $g=.1$. The true Borda gap is exactly $g$: the winner has score $.5+a$, others $.5-a/(K-1)$. Thus this experiment changes $K$ without silently shrinking the Borda gap. Pairwise ties are permitted. There are 100 replications for $K\le10$ and 30 for $K=20,50$; the cap is $10^5$ duels. The larger settings retain the five core modes to limit runtime.

\textbf{B: Borda gap.} At $K=5$ vary the verified gap over $.20,.10,.05,.02,.01$. Use 100 seeds and cap $10^5$ for the first three, 30 seeds and cap $2\cdot10^5$ for the last two. A cap is a censoring threshold, not evidence of successful identification. We do not fit a scaling exponent to saturated points. Decreasing the gap should require substantially more evidence, even when empirical recommendations are often already correct.

\textbf{C: cycles and pairwise heterogeneity.} Add directed circulations of amplitude $.30$ to disjoint triples: increase $p_{ij},p_{j\ell},p_{\ell i}$ and decrease their reverse probabilities. Every row sum is unchanged, so each Borda score and gap is preserved. Matrices remain valid because the base margins and perturbations stay inside $[0,1]$. Compare the resulting $K=3,5,10$ instances against A. This isolates a change in nontransitivity and probability heterogeneity from a change in the objective's difficulty parameter. Include a nearly symmetric three-arm cycle with gap $.01$ and the explicit Condorcet/Borda-disagreement matrix. The construction provides balanced and heterogeneous probabilities; it does not span every practically relevant sparse-comparison structure.

\textbf{D: robustness.} A symmetric $.8$ three-cycle has an exact three-way Borda tie. With 30 seeds and a $10^5$ cap, report certificate failures and recommendations in the maximizing set. Its accuracy statistic is set-valued correctness, not identification of a unique true winner. Near-ties are covered in B and C. No run is removed because it is difficult, wrong, or uncertified.

\subsection{Metrics, uncertainty, and horizon experiments}
The main cost statistic is the \emph{restricted mean} $\E[\min(\tau,H)]$, estimated by the mean charged duel count; every unfinished run contributes its full cap. We retain per-seed counts, quartiles, certification, empirical correctness, allocation matrices, elapsed runtime, and Borda regret. Savings are $100(\bar N_U-\bar N_C)/\bar N_U$; negative values are regressions. These are actual identification costs only for completed runs. If both methods hit $H$, their equal charged cost cannot determine which would certify first with additional samples. A tie at the cap is not evidence of equal efficiency.

Means have 95\% normal standard-error bars in the scaling and regret figures; these are descriptive across replications. Savings intervals use 3,000 paired bootstrap resamples of seed indices with fixed bootstrap seed 879. They quantify simulation variation conditional on the chosen matrix, not uncertainty from sampling the original human population. Fixed-$\delta=.05$ regret curves exploit after certification, otherwise continue learning to the horizon. Separate horizon-aware runs use $T\in\{10^3,10^4,10^5\}$, $\delta=T^{-2}$, the easy $K=3,g=.2$ matrix, and 30 seeds per method per horizon. Those runs illustrate the confidence cost of the theorem's regime; they do not establish its asymptotic rate empirically.

\section{Results: comparisons, allocation, and regret}
\subsection{Total comparison cost and scaling}
Table~\ref{tab:cost} reports certification for \emph{both} methods, including capped and tied settings. The 6.8\% saving on cycle3 contrasts with 31.6\% and 36.7\% extra cost at planted $K=5,10$. Three mechanisms explain the regression. First, each row interval must account for $K-1$ opponents, and simultaneous pair coverage pays a $\log M$ correction for $M=K(K-1)/2$ pairs; the componentwise worst-case average can remain wide. Second, adaptive blocks displace direct Borda samples: only three quarters of COBRA's blocks feed its direct estimator. Third, the added pairwise information need not compensate for those lost direct samples. When it does not certify earlier, COBRA approaches the safeguard's sampling surcharge. None of this contradicts valid interval coverage, which guarantees correctness rather than efficiency.

\begin{table}[!htbp]\centering\small
\begin{tabular}{lrrrrr}\toprule Instance&Uniform duels&COBRA duels&Uniform cert.&COBRA cert.&Saving (\%)\\\midrule
@@TABLE@@
\bottomrule\end{tabular}
\caption{Mean charged total duels and certified runs/total runs for both methods. All cap hits remain included; zero savings at a shared cap cannot rank uncapped identification efficiency. The tie row has no unique target. Savings use Uniform + stop, not fixed-budget Uniform.}\label{tab:cost}\end{table}
\fig{scaling}{.89}{Scaling at fixed true Borda gap $.1$. Error bars are 95\% standard-error intervals. The $10^5$ cap censors difficult settings; a flat capped curve is not a favorable scaling law.}
\fig{gaps}{.89}{Gap sensitivity at five arms. Caps increase to $2\cdot10^5$ for the two smallest gaps. Lines connect the observed restricted means, not a fitted inverse-square law.}

Small gaps in Figure~\ref{fig:gaps} expose a distinction between practical recommendation and certification. An empirical leader can be correct long before its lower confidence bound exceeds every rival's upper bound. At the smallest gaps a large budget can still end without a certificate. This is consistent with the inverse-square gap dependence in the conservative guarantee and existing lower bounds, but cap-limited measurements do not verify that rate. Pairwise margins close to one half no longer trigger a separate diagnostic stage; the remaining difficulty is distinguishing average scores.

\subsection{Matched ablations: allocation and information reuse}
\begin{table}[!htbp]\centering\small
\begin{tabular}{lrrrr}\toprule Method on cycle3&Mean duels&SD&Certified&Accuracy\\\midrule
@@ABTABLE@@
\bottomrule\end{tabular}
\caption{One hundred replications on the three-arm cycle, cap $10^5$. The shadow-allocation controls separate access to reused evidence from pair selection. Accuracy is the empirical recommendation's correctness at stopping or exhaustion.}\label{tab:ablation}\end{table}
\fig{ablation}{.92}{Distribution of total comparisons on cycle3; boxes show median and interquartile range, whiskers use the standard 1.5-IQR rule. Outlier points are omitted visually, but remain in all numeric summaries.}
The comparisons answer three separate questions. \textbf{COBRA versus Random Adaptive Block} tests the weighted-width allocation rule with the same safeguard and confidence certificate. \textbf{COBRA versus Direct-Only Certificate} tests access to shared pairwise confidence information under identical allocation and feedback until the earlier stop. \textbf{COBRA versus Uniform + stop} tests whether the complete method saves total comparisons. Reuse can help against Direct-only while the full method still loses to Uniform; a positive reuse ablation alone is therefore insufficient.

The forward-only control discards a valid implication of each duel: observing $i$ beat $j$ also says $j$ lost to $i$. Its more conservative coverage correction and lost reverse evidence can delay certification. This comparison does not justify counting one physical duel as two independent Borda samples. The no-safeguard candidate tests whether devoting all blocks to uncertain pairs helps enough to offset losing designated samples; its empirical behavior is reported without a termination theorem.

\fig{heatmap}{.94}{Actual unordered-pair counts for a matched-seed, $10^4$-duel cycle5 illustration. Both halves of each symmetric matrix depict the same comparisons. Adaptive allocation differs visibly from random allocation, but concentration by itself does not prove faster identification.}
The heatmaps expose the intended mechanism. Leader/rival edges can receive extra samples, and comparisons unrelated to the current contenders can receive fewer adaptive samples. However, the row-average certificate still includes all opponents. Concentrating comparisons does not make the unobserved contributions vanish, and an early empirical leader can redirect effort away from a future competitor. A variance-aware or gap-difference confidence construction could be less conservative; neither its correctness nor a performance gain is assumed here.

\textbf{Weight sensitivity.} We additionally run 100 matched seeds for each leader--rival weight $w\in\{1,2,4\}$ on cycle3 and the planted five-arm instance, fixing cap $10^5$, $\delta=.05$, $B=128$, checkpoints, and every other parameter. The $w=2$ reruns reproduce all 200 original trajectories exactly (excluding elapsed timing). Table~\ref{tab:weights} reports total comparisons; all settings certify 100/100 correct winners. @@WEIGHTFINDING@@
\begin{table}[!htbp]\centering\small\begin{tabular}{lrrrr}\toprule
Instance&$w=1$&$w=2$&$w=4$&Certified per weight\\\midrule
@@WEIGHTS@@
\bottomrule\end{tabular}\caption{Weight sensitivity: mean total duels over 100 fixed seeds per cell. Other incident-pair weights remain 1. The default is $w=2$.}\label{tab:weights}\end{table}

\subsection{Accuracy, regret, and computation}
\fig{regret}{.87}{Strong Borda regret on cycle3 at fixed error level $.05$, comparing the three main certificate-based methods. Correct certification permits zero-regret self-duels. Uncertified runs continue exploring and remain in the mean.}
Confidence-based stopping reduces unnecessary exploration relative to a baseline forced to continue to its cap. We therefore use Uniform + stop as the substantive denominator for comparison savings. On cycle3 at $10^5$ rounds, mean regret is about 1,011 for Uniform + stop versus 884 for COBRA. Empirical accuracy can be high even where certification is incomplete: a high success rate from finitely many seeds does not replace a confidence certificate. Regret depends on which pairs are explored as well as on the stopping count.

The separate horizon-aware experiment uses tighter confidence as $T$ increases, asking for stronger evidence before exploitation. Three finite horizons cannot demonstrate an asymptotic logarithmic law; that statement rests on the proof and its assumptions. Supporting accuracy-by-budget and horizon-aware plots are retained with the study outputs. The main regret figure deliberately shows only the common fixed-confidence regime.

The primary execution took @@MINUTES@@ minutes in the recorded Python environment. On the ten-arm planted setting, mean elapsed times were @@CPUU@@ seconds/run for Uniform + stop and @@CPUC@@ for COBRA, with different duel counts. Supporting per-duel runtime plots separate these effects. Matrix scans add computation as predicted by their $O(K^2)$ cost; compute savings are not established.

\section{Implementation on real preference datasets}
\textbf{Key distinction: on each full real dataset, both methods achieve 30/30 correct empirical recommendations but 0/30 certified identifications within 500,000 duels.} Correct empirical recommendations are not statistically certified identifications. Since all these runs hit the cap, the experiments cannot determine which method would certify first if given additional samples.
\subsection{World Junior Ladies Free Skate: full 24-alternative profile}
We use PrefLib's Skate collection~\cite{preflib,skate}, file \texttt{00006-00000028.soc}, titled \emph{Worldjunior Ladies Free Skate}. It contains nine complete strict judge rankings over 24 alternatives. For each pair, $p_{ij}$ is the fraction of judges placing $i$ ahead of $j$. No smoothing or restriction to a three-arm cycle is used. The full profile has two strict three-cycles but also a strict Condorcet winner. Cycles among other alternatives do not rule out a global Condorcet winner. This is a human-judgment application with contradictory local preferences, not an example of full-profile Condorcet nonexistence.

Candidate 24 uniquely maximizes Borda, with gap $.0242$. This previously used profile was fixed before the new performance runs; here we evaluate its full matrix. We make no population-level claim from nine judges. The cap-limited costs and certificate counts appear in Table~\ref{tab:cost}.

\subsection{Electoral Reform Society: full 26-candidate election}
We use PrefLib's ERS collection~\cite{preflib,ers}, file \texttt{00007-00000005.soi}, titled \emph{ERS Set 5}. It has 104 ballots over 26 candidates, with incomplete strict rankings. Ranked candidates precede unranked candidates; all unranked candidates on a ballot are tied jointly last. A tied pair contributes $1/2$ to each candidate. This completion rule is explicit and affects the resulting Borda scores. We use every candidate and ballot, with no smoothing. Pairwise ties at probability one half are valid inputs.

The full matrix has 13 strict three-cycles and no strict Condorcet winner. We selected it before performance testing as the first previously audited ERS profile with those properties and a unique Borda maximum. Borda remains well-defined despite the absence of an alternative that beats everyone.

\begin{table}[!htbp]\centering\small\begin{tabular}{lrrrrr}\toprule
Profile&$K$&Voters&3-cycles&Borda ID&Gap\\\midrule
@@REALTABLE@@
\bottomrule\end{tabular}\caption{Full empirical profiles. Candidate IDs are the original one-based identifiers. Cycles count unordered triples satisfying a strict directed majority cycle.}\label{tab:real}\end{table}

For both datasets, we simulate independent Bernoulli duels from the resulting empirical matrix. These are not new live human experiments: the matrix is treated as ground truth conditional on the observed rankings. For ERS, the half-vote treatment of tied preferences corresponds to a fair random tie outcome. Fresh simulated draws do not increase the number of independent human voters. The original rankings, source URLs, processing rule, derived matrices, and matrix hashes are retained so the numerical evaluation can be audited.

\begin{table}[!htbp]\centering\small\begin{tabular}{lrrr}\toprule
Full profile and method&Certified&Accuracy&Seconds/run\\\midrule
@@PERFORMANCE@@
\bottomrule\end{tabular}\caption{Real-data recommendations and certificates at the common $5\cdot10^5$ cap. Accuracy is empirical Borda-target correctness, not a certificate.}\end{table}
\fig{real_regret}{.88}{Full-profile Borda regret from actual requested pairs. Solid lines denote skating and dashed lines ERS. An adaptive method can have different regret even when both methods exhaust the same comparison cap.}
At $T=500{,}000$, mean skating regret is about 245,252 for Uniform + stop versus 203,748 for COBRA; ERS gives 58,367 versus 51,246. These lower-regret exploration trajectories are a positive observation even though neither method certifies either profile. They do not establish lower identification cost or an asymptotic regret improvement.

\section{Discussion, limitations, and reflection}
\begin{table}[!htbp]\centering\small\begin{tabular}{lrr}\toprule
Instance&Saving (\%)&95\% paired bootstrap interval\\\midrule
@@CITABLE@@
\bottomrule\end{tabular}\caption{Simulation uncertainty for restricted-mean savings. Negative intervals indicate a regression under the tested matrix and cap; they are not uncertainty intervals over new human populations.}\label{tab:ci}\end{table}

\textbf{What does the comparison establish?} The matched cycle3 ablation separates reused evidence from its acquisition cost; scaling tests whether the benefit survives larger $K$ at fixed gap. A favorable reuse ablation is insufficient to recommend replacing uniform learning.

\textbf{Why might uniform learning be enough?} A designated observation estimates a whole Borda score directly. Pairwise row intervals average worst-case component endpoints and discard cancellation. If they fail to compensate for lost direct samples, the three-safe/one-adaptive schedule approaches a $4/3$ sampling surcharge. Valid coverage does not imply an efficient allocation.

\textbf{What changes without a Condorcet winner?} The active code does not test its existence. Cycles, the full ERS election, and the disagreement matrix all use Borda; the disagreement target is $B$. The archived project's hierarchical objective differs, so its results are not a controlled numerical baseline here.

\textbf{What is guaranteed?} Certificates are simultaneously valid under stationary independent draws. The safeguard gives a gap-dependent completion bound and finite-horizon regret guarantee; the adaptive heuristic has no sharper theorem. Exact ties invalidate strict-winner termination, and small gaps may exceed the cap. Neither gap-free $O(K\log T)$ regret nor universal $\sqrt K$ scaling follows.

\textbf{What remains uncertain?} Planted matrices with tied suboptimal scores may understate benefits from eliminating inferior arms. Two real profiles and 30 large-instance seeds limit generalization; incomplete-ballot completion is a modeling choice. Bootstrap intervals are exploratory and unadjusted for multiple instances. We do not establish superiority to SECS or other active-ranking methods. Difference-specific confidence bounds and successive elimination are useful next comparisons under the same objective and budget.

\textbf{Validation and learning.} The 21 tests cover the original matrix, confidence, certificate, regret, tie, and reproducibility checks plus the actual block schedule, shortened blocks, $N_0$, and weight defaults. The policy API has no true-matrix or gap parameter. Ground truth is confined to construction, simulation, and evaluation. Numerical summaries are regenerated from per-seed outputs; replay checks reproduce selected original trajectories and all default-weight sensitivity trajectories. Code, pinned dependencies, seeds, configurations, raw results, plotting scripts, tests, and final LaTeX/PDF are retained in the repository. The probability guarantee rests on the proof, not on an absence of observed errors.

\section{Conclusion}
Does COBRA consistently outperform uniform Borda? \textbf{No.} Its advantage depends on preference structure, confidence-bound conservatism, and sampling allocation. The 6.8\% cycle3 saving accompanies 31.6\% and 36.7\% extra cost at planted $K=5,10$. Reused evidence may not repay displaced direct samples. Both real profiles remain uncertified at the cap; their uncapped efficiency ordering is unknown. The uniform safeguard preserves a conservative guarantee, but Theorem 4 needs a supplied finite horizon with $\delta=T^{-2}$. It establishes neither anytime regret nor adaptive improvement in gap dependence.

\FloatBarrier
\addcontentsline{toc}{section}{References}
\bibliographystyle{plain}\bibliography{references}
\noindent\textbf{Project repository:} \url{https://github.com/pranav1785/COBRA}
\end{document}
